\chapter{System Design}
\label{ch:design}

This chapter presents the design of \textsc{EdgeGuard} and argues the properties claimed for
each component. \Cref{sec:design:overview} gives the architecture. \Cref{sec:design:merkle}
covers weight attestation, \cref{sec:design:monitor} the activation monitor,
\cref{sec:design:classifier} the threat classifier, \cref{sec:design:llm} the
language-model-specific scores, \cref{sec:design:explain} the explanations and
\cref{sec:design:adapter} the runtime adapters.

% ---------------------------------------------------------------------------
\section{Architecture}
\label{sec:design:overview}

\begin{figure}[t]
\centering
\begin{tikzpicture}[
  blk/.style={draw=egblue, fill=eglight, rounded corners=2pt, align=center, font=\small, minimum height=1.0cm, minimum width=3.4cm},
  io/.style={draw=eggray, rounded corners=2pt, align=center, font=\small, minimum height=0.8cm, minimum width=2.6cm},
  outb/.style={draw=egorange, fill=egorange!8, rounded corners=2pt, align=center, font=\small, minimum height=1.0cm, minimum width=3.2cm},
  arr/.style={-{Stealth}, eggray, thick}]
\node[io] (w) at (0,0) {Weight tensors\\$\{L_i\}$ in DRAM};
\node[io] (a) at (0,-2.6) {Per-layer INT8\\activations $\{\mathbf{a}_\ell\}$};
\node[blk] (att) at (4.4,0) {\textbf{Attestation}\\Merkle verify + ECDSA};
\node[blk] (mon) at (4.4,-2.6) {\textbf{Activation monitor}\\$Z_\ell$, $D_\ell$ per layer};
\node[blk] (llm) at (4.4,-4.9) {\textbf{LLM scores}\\memorisation, entropy};
\node[blk, minimum height=6.2cm, minimum width=2.6cm] (clf) at (8.8,-2.45) {\textbf{Threat}\\\textbf{classifier}\\priority\\rules};
\node[outb] (res) at (12.4,-2.45) {Label, confidence,\\severity,\\explanation};
\node[io, font=\scriptsize, minimum width=2.2cm] (ref) at (4.4,1.6) {signed root $(r,\sigma)$};
\node[io, font=\scriptsize, minimum width=2.2cm] (cal) at (1.2,-4.0) {calibration\\$\boldsymbol\mu,\hat\Sigma$};
\draw[arr] (w) -- (att);
\draw[arr] (a) -- (mon);
\draw[arr] (a.south) |- (llm.west);
\draw[arr] (ref) -- (att);
\draw[arr] (cal.east) -- ++(0.5,0) |- (mon.south west);
\draw[arr] (att.east) -- node[above, font=\scriptsize]{$\mathcal{M}$} (att.east -| clf.west);
\draw[arr] (mon.east) -- node[above, font=\scriptsize]{$\mathcal{A}$} (mon.east -| clf.west);
\draw[arr] (llm.east) -- node[above, font=\scriptsize]{$s_{\text{mem}},H$} (llm.east -| clf.west);
\draw[arr] (clf) -- (res);
\end{tikzpicture}
\caption{\textsc{EdgeGuard} runtime architecture. The two detectors are independent: attestation
reads only weights; the monitor reads only activations. $\mathcal{M}$: tensors failing
attestation; $\mathcal{A}$: anomalous layers; $s_{\text{mem}},H$: memorisation and entropy
scores. The classifier fuses these outputs.}
\label{fig:architecture}
\end{figure}

\Cref{fig:architecture} shows the runtime structure. Two detectors operate on disjoint inputs.
Attestation reads only the weights and compares them with a signed reference. The activation
monitor reads only the activations of the current inference and compares them with calibration
statistics. A third, model-specific stage computes two scalar scores from late-layer
activations. The classifier fuses the three outputs by priority rules into a single label with a
confidence, a severity and an explanation.

The separation is deliberate. Because the detectors share no inputs, an attacker who controls
one input channel, say the prompt, cannot use it to influence the other detector's evidence.
This is the basis of the complementarity argument in \cref{sec:design:complement}.

The offline phase has two steps. The vendor builds and signs the Merkle tree
(\cref{alg:build}), and the deployer calibrates the monitor on clean inputs
(\cref{alg:calibrate}). Both produce small artefacts: a tree of $2n$ digests with one
signature, and per-layer means and covariances.

% ---------------------------------------------------------------------------
\section{Weight attestation}
\label{sec:design:merkle}

\subsection{Leaf construction}

Each weight tensor $L_i$ is hashed together with its shape:
\begin{equation}
  h_i = \Hash\big(\mathrm{enc}(\mathrm{shape}(L_i)) \,\|\, \mathrm{bytes}(L_i)\big),
  \label{eq:leaf}
\end{equation}
where $\mathrm{enc}$ writes each dimension as a big-endian 64-bit integer and
$\mathrm{bytes}$ is the row-major (C-order) byte serialisation. Including the shape means that
two tensors with the same bytes but different shapes, for example a $64\times16$ and a
$16\times64$ matrix, produce different leaves. Leaves are ordered by the model's tensor order,
which for HuggingFace-style checkpoints is deterministic.

\subsection{Tree construction and signing}

\begin{algorithm}[t]
\caption{\textsc{BuildAttestation}$(\{L_i\}_{i=0}^{n-1}, sk)$}
\label{alg:build}
\begin{algorithmic}[1]
\State $\mathit{level} \gets [\,h_0,\dots,h_{n-1}\,]$ \Comment{leaves per \cref{eq:leaf}}
\While{$|\mathit{level}| > 1$}
  \If{$|\mathit{level}|$ is odd} append a copy of the last element \EndIf
  \State $\mathit{level} \gets [\,\Hash(\mathit{level}[2j] \,\|\, \mathit{level}[2j{+}1]) : j = 0,\dots \,]$
\EndWhile
\State $r \gets \mathit{level}[0]$
\State $\sigma \gets \textsc{ECDSA-Sign}_{P\text{-}256}(sk, r)$
\State \Return $(r, \sigma, [h_0,\dots,h_{n-1}], n)$
\end{algorithmic}
\end{algorithm}

\Cref{alg:build} builds the tree bottom-up, duplicating the last node of any odd level, and
signs the root. The stored artefact consists of the root, the signature, the leaf digests, the
leaf count $n$ and the public key.

\begin{figure}[t]
\centering
\begin{tikzpicture}[
  nd/.style={draw=eggray, rounded corners=1.5pt, minimum width=1.05cm, minimum height=0.55cm, font=\scriptsize},
  bad/.style={nd, draw=egorange, fill=egorange!12, thick},
  sib/.style={nd, draw=egblue, fill=eglight},
  level distance=1.05cm]
\foreach \i in {0,...,7} {
  \ifnum\i=5 \node[bad] (l\i) at (\i*1.45,0) {$h_{\i}$};
  \else\ifnum\i=4 \node[sib] (l\i) at (\i*1.45,0) {$h_{\i}$};
  \else \node[nd] (l\i) at (\i*1.45,0) {$h_{\i}$}; \fi\fi
  \node[font=\scriptsize, text=eggray] at (\i*1.45,-0.6) {$L_{\i}$};
}
\node[nd] (p0) at (0.725,1.1) {$h_{01}$};
\node[nd] (p1) at (3.625,1.1) {$h_{23}$};
\node[bad] (p2) at (6.525,1.1) {$h_{45}$};
\node[sib] (p3) at (9.425,1.1) {$h_{67}$};
\node[sib] (q0) at (2.175,2.2) {$h_{0..3}$};
\node[bad] (q1) at (7.975,2.2) {$h_{4..7}$};
\node[bad] (r) at (5.075,3.3) {$r$};
\foreach \a/\b in {l0/p0,l1/p0,l2/p1,l3/p1,l4/p2,l5/p2,l6/p3,l7/p3,p0/q0,p1/q0,p2/q1,p3/q1,q0/r,q1/r}
  \draw[eggray] (\a.north) -- (\b.south);
\node[font=\scriptsize, align=left, anchor=west] at (11.0,2.2) {\textcolor{egorange!80!black}{$\blacksquare$} path changed by\\ \quad a flip in $L_5$\\[2pt]
\textcolor{egblue}{$\blacksquare$} siblings in the\\ \quad proof for $L_5$};
\end{tikzpicture}
\caption{Merkle tree over eight weight tensors. A change in $L_5$ alters every node on its path
to the root (orange). The membership proof for $L_5$ consists of the three sibling digests
(blue), so a single tensor can be checked against the signed root with
$\lceil\log_2 n\rceil$ extra hash evaluations.}
\label{fig:merkle}
\end{figure}

\subsection{Verification and localisation}

\begin{algorithm}[t]
\caption{\textsc{Verify}$(\{L'_i\}, (r,\sigma,[h_i],n), pk)$}
\label{alg:verify}
\begin{algorithmic}[1]
\Require $|\{L'_i\}| = n$ \Comment{leaf count is fixed by the reference}
\State $\mathit{sig\_ok} \gets \textsc{ECDSA-Verify}(pk, r, \sigma)$
\State $h'_i \gets$ leaf digest of $L'_i$ per \cref{eq:leaf}, for all $i$
\State $r' \gets$ root of the tree over $[h'_0,\dots,h'_{n-1}]$ \Comment{\cref{alg:build}, lines 1--5}
\If{$r' = r$} \Return $(\mathit{valid}=\mathit{sig\_ok}, \mathcal{M}=\emptyset)$ \EndIf
\State $\mathcal{M} \gets \{\, i : h'_i \neq h_i \,\}$
\State \Return $(\mathit{valid}=\textsf{false}, \mathcal{M})$
\end{algorithmic}
\end{algorithm}

\Cref{alg:verify} recomputes every leaf, rebuilds the root and compares it with the signed root.
If they differ, it compares leaves to find the modified tensors. The signature check confirms
that the \emph{reference} is authentic. The root comparison confirms that the \emph{current
weights} match the reference. Both must pass.

\begin{proposition}[Integrity, G1]\label{prop:integrity}
If the stored reference is authentic and $L'_i \neq L_i$ for some $i$ (in bytes or shape),
\cref{alg:verify} returns $\mathit{valid}=\textsf{false}$, unless the attacker has found a
SHA-256 second preimage.
\end{proposition}
\begin{proof}[Argument]
If $L'_i\neq L_i$, the inputs to \cref{eq:leaf} differ, so $h'_i = h_i$ would be a second
preimage. Hence $h'_i\neq h_i$. Now suppose $r'=r$. Walking up the path from leaf $i$, there is a
first level at which the recomputed node equals the reference node while a child differs. At
that node, two different 64-byte inputs have the same SHA-256 digest, which is a collision.
Hence $r'\neq r$ and the algorithm returns \textsf{false}.
\end{proof}

\begin{proposition}[Exact localisation, G2]\label{prop:localise}
Under the same assumption, the set $\mathcal{M}$ returned by \cref{alg:verify} equals
$\{i : L'_i \neq L_i\}$.
\end{proposition}
\begin{proof}[Argument]
$\mathcal{M}$ is computed by direct comparison of each fresh leaf with its stored counterpart.
By the first step of the previous argument, $h'_i\neq h_i$ if and only if $L'_i\neq L_i$, unless
a second preimage exists.
\end{proof}

\begin{remark}[Duplicate-leaf padding]\label{rem:padding}
Padding an odd level with a copy of its last node means that the leaf lists $[h_0,\dots,h_{k}]$
and $[h_0,\dots,h_{k},h_{k}]$ can yield the same root. This is the root of a known weakness in
some Merkle-tree implementations. \Cref{alg:verify} is not affected because it requires the
fresh leaf count to equal the signed $n$ and it compares leaves individually. Any variant that
accepts a variable number of tensors must, however, either sign $n$ together with $r$ or use
distinct domain-separation prefixes for leaves, internal nodes and padding. The current
implementation relies on the fixed count. Adding domain separation is listed as future work.
\end{remark}

\paragraph{Cost.} Verification is dominated by hashing every byte of every tensor, so its cost is
$\Theta(\sum_i |L_i|)$ regardless of the tree. The tree adds $n-1$ internal hashes, which are
negligible, and leaf comparison costs $O(n)$ digest comparisons. The value of the tree structure is
therefore \emph{not} faster full verification. Its value is that any single tensor can be checked
against the signed root with a proof of $\lceil\log_2 n\rceil$ digests, without re-hashing the
others. This enables the amortised schemes in \cref{sec:disc:amortise}, which are needed because
full verification of language-model-sized weights is too slow to run on every inference
(\cref{sec:eval:merkle}). The current implementation generates inclusion proofs (\code{MerkleTree.get\_proof}) but does
not yet include a stand-alone proof verifier or a partial-verification mode. Both are planned
(\cref{ch:conclusion}).

% ---------------------------------------------------------------------------
\section{INT8 activation monitor}
\label{sec:design:monitor}

\subsection{Calibration}

\begin{algorithm}[t]
\caption{\textsc{Calibrate}$(\{\mathbf{a}_{\ell,t}\}, d_{\max})$}
\label{alg:calibrate}
\begin{algorithmic}[1]
\For{each monitored layer $\ell$}
  \State $t \gets 0$; $\boldsymbol\mu \gets \mathbf{0}$; $M \gets \mathbf{0}$
  \For{each clean sample $\mathbf{a}_{\ell,t}$} \Comment{INT8 values, cast to float64 for accumulation}
    \State $t \gets t+1$; $\boldsymbol\delta \gets \mathbf{a}_{\ell,t} - \boldsymbol\mu$; $\boldsymbol\mu \gets \boldsymbol\mu + \boldsymbol\delta/t$
    \State $M \gets M + \boldsymbol\delta(\mathbf{a}_{\ell,t}-\boldsymbol\mu)^\top$ \Comment{diagonal only if $d_\ell > d_{\max}$}
  \EndFor
  \State $\hat\Sigma_\ell \gets M/(t-1)$; $\boldsymbol\sigma_\ell \gets \sqrt{\mathrm{diag}(\hat\Sigma_\ell)}$
  \State $\hat\Sigma_\ell^{+} \gets \textsc{PseudoInverse}(\hat\Sigma_\ell)$ if $d_\ell \le d_{\max}$
\EndFor
\end{algorithmic}
\end{algorithm}

\Cref{alg:calibrate} applies Welford's update \cref{eq:welford} to each monitored layer. The
activations are the INT8 values the runtime produces, so the statistics are in the same units
that are seen at inference time. For layers wider than $d_{\max}$ (default 512), only the
diagonal is kept, because a dense covariance costs $O(d^2)$ memory: 1\,MiB for $d=512$ at double
precision, but 128\,MiB for $d=4096$. The pseudo-inverse uses an eigendecomposition for symmetric
matrices (\code{scipy.linalg.pinvh}), which tolerates the singular covariances that arise when
some channels are constant or perfectly correlated.

\subsection{Scoring and decision}

For each monitored layer $\ell$ and test activation $\mathbf{a}_\ell$, the monitor computes
\begin{equation}
  Z_\ell = \max_j \frac{|a_{\ell,j}-\mu_{\ell,j}|}{\sigma_{\ell,j}+\epsilon},
  \qquad
  D_\ell = \sqrt{(\mathbf{a}_\ell-\boldsymbol\mu_\ell)^\top\hat\Sigma_\ell^{+}(\mathbf{a}_\ell-\boldsymbol\mu_\ell)},
  \label{eq:scores}
\end{equation}
with $D_\ell$ falling back to $Z_\ell$ when no covariance is stored. A layer is \emph{anomalous}
if $Z_\ell>\tau_z$ \emph{or} $D_\ell>\tau_D$. The input is anomalous if any layer is. The
implementation also reports a weighted combination of normalised $Z_\ell$ and $D_\ell$ for use in
explanations, but the decision uses the two thresholds directly.

\subsection{Choosing the thresholds}
\label{sec:design:threshold}

The $z$ statistic is a maximum over $d_\ell$ channels, and the input-level decision is a further
maximum over layers. The effective number of comparisons is therefore $m=\sum_\ell d_\ell$, and a
per-channel threshold that would be conservative in isolation produces a substantial input-level
false-positive rate by \cref{eq:fwer}. The design offers two ways to set $\tau_z$:

\begin{enumerate}[leftmargin=*]
  \item \textbf{Bonferroni:} $\tau_z = z^*(\alpha)$ from \cref{eq:bonf}, for a target input-level
        false-positive rate $\alpha$. This needs only $m$.
  \item \textbf{Empirical:} $\tau_z$ is the $(1-\alpha)$ quantile of the input-level maximum
        $\max_\ell Z_\ell$ on a held-out clean set. This adapts to correlation and non-Gaussian
        tails but needs enough held-out data to estimate a tail quantile.
\end{enumerate}

\Cref{sec:eval:fpr} shows that the Bonferroni threshold achieves close to its target on the
evaluation data. The fixed default of $\tau_z=4$ used in the main experiments was chosen before
this analysis and is retained so that its effect can be reported.

% ---------------------------------------------------------------------------
\section{Threat classifier}
\label{sec:design:classifier}

\begin{figure}[t]
\centering
\begin{tikzpicture}[
  dec/.style={draw=eggray, diamond, aspect=2.4, align=center, font=\scriptsize, inner sep=1pt},
  lab/.style={draw=egorange, fill=egorange!8, rounded corners=2pt, align=center, font=\scriptsize, minimum width=2.4cm, minimum height=0.6cm},
  arr/.style={-{Stealth}, eggray},
  node distance=0.5cm]
\node[dec] (d1) {$\mathcal{M}\neq\emptyset$?};
\node[dec, below=0.55cm of d1] (d2) {$s_{\text{mem}}>\theta_m$?};
\node[dec, below=0.95cm of d2] (d3) {$H<\theta_H$?};
\node[dec, below=0.95cm of d3] (d4) {$\mathcal{A}\neq\emptyset$?};
\node[dec, right=1.2cm of d1] (d1b) {$|\mathcal{M}|>0.3n$?};
\node[lab, right=0.9cm of d1b, yshift=0.45cm] (sc) {SUPPLY\_CHAIN\_BACKDOOR};
\node[lab, right=0.9cm of d1b, yshift=-0.45cm] (wt) {WEIGHT\_TAMPER};
\node[lab, right=4.4cm of d2] (ex) {TRAINING\_DATA\_EXFILTRATION};
\node[dec, right=1.2cm of d3] (d3b) {base label\\$=$ ADV?};
\node[lab, right=0.9cm of d3b, yshift=0.45cm] (jb) {JAILBREAK};
\node[lab, right=0.9cm of d3b, yshift=-0.45cm] (pi) {PROMPT\_INJECTION};
\node[dec, right=1.2cm of d4] (d4b) {early share\\$\ge 0.5$?};
\node[lab, right=0.9cm of d4b, yshift=0.45cm] (adv) {ADVERSARIAL\_INPUT};
\node[lab, right=0.9cm of d4b, yshift=-0.45cm] (tf) {TRANSIENT\_FAULT};
\node[lab, below=0.55cm of d4] (cl) {CLEAN};
\draw[arr] (d1) -- node[left, font=\tiny]{no} (d2);
\draw[arr] (d2) -- node[left, font=\tiny]{no} (d3);
\draw[arr] (d3) -- node[left, font=\tiny]{no} (d4);
\draw[arr] (d4) -- node[left, font=\tiny]{no} (cl);
\draw[arr] (d1) -- node[above, font=\tiny]{yes} (d1b);
\draw[arr] (d1b.east) -- node[above, font=\tiny]{yes} (sc.west);
\draw[arr] (d1b.east) -- node[below, font=\tiny]{no} (wt.west);
\draw[arr] (d2) -- node[above, font=\tiny]{yes} (ex);
\draw[arr] (d3) -- node[above, font=\tiny]{yes} (d3b);
\draw[arr] (d3b.east) -- node[above, font=\tiny]{yes} (jb.west);
\draw[arr] (d3b.east) -- node[below, font=\tiny]{no} (pi.west);
\draw[arr] (d4) -- node[above, font=\tiny]{yes} (d4b);
\draw[arr] (d4b.east) -- node[above, font=\tiny]{yes} (adv.west);
\draw[arr] (d4b.east) -- node[below, font=\tiny]{no} (tf.west);
\end{tikzpicture}
\caption{Priority decision procedure of the language-model threat classifier. $\mathcal{M}$:
tensors failing attestation; $s_{\text{mem}}$: memorisation score; $H$: late-layer entropy;
$\mathcal{A}$: anomalous layers; ``early share'': fraction of anomalous layers lying in the first
40\% of the model.}
\label{fig:classifier}
\end{figure}

\Cref{fig:classifier} gives the decision procedure. Its order encodes a judgement about
evidence strength. A hash mismatch is cryptographic evidence and takes precedence over
everything. The model-specific scores come next. Generic activation anomalies come last, and
their placement within the network decides between an input-borne attack and a fault.

\paragraph{Weight evidence.} Any attestation failure produces WEIGHT\_TAMPER, upgraded to
SUPPLY\_CHAIN\_BACKDOOR when more than 30\% of tensors are affected. The rationale is that a
runtime bit-flip attack touches a handful of bits, whereas a replaced or re-poisoned checkpoint
changes many tensors. The heuristic is inherently weak. A supply-chain attacker who modifies a
single tensor is indistinguishable, from the weights alone, from a runtime flip in that tensor.
\Cref{sec:eval:labels} shows this happening.

\paragraph{Activation evidence.} With intact weights, an anomaly concentrated in early layers is
labelled an adversarial input, on the reasoning that crafted inputs perturb the earliest
representations. An anomaly concentrated in late layers is labelled a transient fault, for
example a soft error in on-chip SRAM. The split uses the fraction of \emph{anomalous} layers
that are early, compared against 0.5, with ``early'' meaning the first 40\% of layers. One
consequence is worth noting before the evaluation. If every layer is anomalous, the early share
is 0.4 and the label is TRANSIENT\_FAULT, even for an input attack. A rule that compared the
\emph{rate} of anomalies among early layers with the rate among late layers would not have this
bias. The issue is examined with data in \cref{sec:eval:labels}.

\paragraph{Confidence and severity.} Each label carries a heuristic confidence: 0.90 plus a small
increase with the fraction of tampered tensors for weight tampering, a value rising with the
early share and the peak $z$ for adversarial inputs, and so on. These are ranking aids for
operators, not calibrated probabilities. Severity is a fixed mapping from label to
\{none, low, medium, high, critical\}, downgraded to medium when confidence is below 0.5.

% ---------------------------------------------------------------------------
\section{Language-model-specific scores}
\label{sec:design:llm}

Two scalar scores are computed over the late layers $\mathcal{L}$ (the last 60\% of
tensors in model order). Both are defined exactly as implemented.

\paragraph{Memorisation score.} For each late layer, with mean absolute activation $\bar a$ and
standard deviation $s$,
\begin{equation}
  s_\ell = \min\!\Big(1,\; \frac{\bar a}{50}\cdot\frac{1}{1+s/5}\Big),
  \qquad s_{\text{mem}} = \frac{1}{|\mathcal{L}|}\sum_{\ell\in\mathcal{L}} s_\ell ,
  \label{eq:mem}
\end{equation}
with $s_\ell=1$ if $s=0$. The score is high when late activations are large and nearly
constant, which is the signature the extraction-probe simulator assumes. The constants 50 and 5
are scale choices for INT8 activations near the calibration distribution used here (mean 10,
standard deviation 5). They are not derived from a real model and would need re-fitting.

\paragraph{Late-layer entropy.} For each late layer, a softmax with a data-dependent temperature
is applied to the activation vector,
\begin{equation}
  p_j \propto \exp\!\Big(\frac{a_j - \max_k a_k}{\max(\overline{|a - \max a|},\,1)}\Big),
  \qquad
  H_\ell = -\frac{\sum_j p_j\log p_j}{\log d},
  \label{eq:entropy}
\end{equation}
and $H$ is the mean of $H_\ell$ over late layers. Values near 1 indicate diffuse activations;
values near 0 indicate a few dominant channels. Because the temperature is the mean deviation
from the maximum, the score is invariant to adding a constant to every channel and to
multiplying every channel by a positive constant. Shifting or uniformly narrowing a layer's
activations therefore cannot lower it, which matters in \cref{sec:eval:labels}.

\paragraph{Thresholds.} The evaluation uses $\theta_m=0.55$ and $\theta_H=0.30$.

\paragraph{Unreachable labels.} The label set defines WEIGHT\_THEFT\_PROBE and UNKNOWN, but the
decision procedure in \cref{fig:classifier} never returns them. They are kept as placeholders for
future rules. The classifier therefore has eight reachable labels, and all results are reported
over those eight.

% ---------------------------------------------------------------------------
\section{Complementarity of the two detectors}
\label{sec:design:complement}

The design argument for combining attestation with activation monitoring can be stated
precisely. Let $E_W$ be the event that weights are modified, and $E_X$ the event that the
current input is adversarial. Attestation detects $E_W$ with certainty under
\cref{prop:integrity}, independently of the input. The monitor detects an input when its
activations leave the calibration region, whatever the cause. A synergy attacker realises $E_W$
and chooses the input to keep the activations inside the calibration region. Against the monitor
alone, such an attacker can succeed by construction if the input space is rich enough. Against the
combination, the attacker must defeat attestation, which requires a hash collision or a
compromise of the verifier, both excluded by the threat model. The combination therefore turns a
statistical game that the attacker can win into a cryptographic one that the attacker cannot.

The argument has a limit. It holds only at verification time. A flip made after verification
and before the next one is invisible to attestation during that window (\cref{sec:disc:toctou}).

% ---------------------------------------------------------------------------
\section{Explanations}
\label{sec:design:explain}

Each result carries a plain-text explanation assembled from the same quantities used for the
decision. Nothing is generated by a separate model. The following explanation was produced for a
simulated sign-bit flip in the tiny profile. It is reproduced from the implementation's output,
with a check-mark glyph before ``valid'' omitted:

\begin{lstlisting}[language={},numbers=none,frame=single]
LLM THREAT: WEIGHT_TAMPER
  Base classification: WEIGHT_TAMPER (conf=0.91)
  Memorization probe score: 0.099 (threshold=0.55)
  Late-layer entropy: 0.965 (threshold=0.3)
  Tampered layers: ['model.layers.0.mlp.gate_proj.weight']
  Merkle root signature: valid
\end{lstlisting}

The last line reports that the \emph{reference} root's signature verified, not that the weights
are intact. The wording could mislead an operator, and changing it to ``reference signature
valid; current weights do NOT match'' is a small usability fix recorded in
\cref{ch:conclusion}. For an activation-only detection, the explanation lists the anomalous
layers with their $Z_\ell$ and $D_\ell$ values next to the thresholds, so that an analyst can
see how far each layer exceeded its threshold.

% ---------------------------------------------------------------------------
\section{Runtime adapters}
\label{sec:design:adapter}

\paragraph{PyTorch.} For development, activations are captured with forward hooks registered on
named modules. A context manager installs and removes the hooks and collects one activation
tensor per module per call.

\paragraph{QNN.} For an NPU deployment, the adapter is defined against a minimal protocol rather
than the vendor SDK, so that it can be tested without hardware:

\begin{lstlisting}
class QNNGraph(Protocol):
    def execute(self, inputs: dict[str, np.ndarray]) -> dict[str, np.ndarray]: ...
    def get_layer_names(self) -> list[str]: ...
    def get_intermediate_outputs(self, layer: str) -> np.ndarray: ...
\end{lstlisting}

A configuration maps each monitored \textsc{EdgeGuard} layer name to one or more graph node
names, and an option clips values to the INT8 range. A concrete implementation over the QNN C
API would expose the chosen intermediate tensors as additional graph outputs when the context
binary is built. Writing that implementation and measuring its overhead on a device is the main
engineering task of the next phase.
